% 第四章 技术实现
\section{技术实现}

\subsection{密码学核心模块}

密码学核心模块（shamir.ts）是整个系统的基础设施，提供以下关键功能：

\subsubsection{双多项式份额生成}

标准 Shamir 方案仅拆分份额值，无法在不暴露值的情况下验证份额正确性。本系统引入第二个\textbf{盲因子多项式} $Q(x)$，构成双多项式结构：

\[
\begin{cases}
P(x) = s + a_1x + a_2x^2 + \cdots + a_{t-1}x^{t-1} \pmod{p} \\
Q(x) = R + b_1x + b_2x^2 + \cdots + b_{t-1}x^{t-1} \pmod{p}
\end{cases}
\]

其中 $P(0) = s$（主密钥），$Q(0) = R$（随机主盲因子）。份额 $i$ 包含值 $v_i = P(i)$ 和盲因子 $r_i = Q(i)$。对应的 Pedersen 承诺为 $C_i = r_iG + v_iH$。

\vspace{0.5em}
\noindent 核心代码实现如下：

\begin{lstlisting}[language=TypeScript, caption=双多项式份额生成（shamir.ts）]
generatePedersenShares(masterKey: string, n: number, t: number) {
    const secretNum = BigInt('0x' + masterKey)
    const pCoeffs: bigint[] = [secretNum]     // P(x): P(0) = masterKey
    const qCoeffs: bigint[] = [randomCoeff()] // Q(x): Q(0) = random R

    for (let i = 1; i < t; i++) {
        pCoeffs.push(generateRandomCoefficient())
        qCoeffs.push(generateRandomCoefficient())
    }

    for (let i = 1; i <= n; i++) {
        const x = BigInt(i)
        shares.push({
            index: i,
            value: evaluatePoly(pCoeffs, x).toString(16),     // v_i = P(i)
            blindingFactor: evaluatePoly(qCoeffs, x).toString(16), // r_i = Q(i)
        })
    }
    return { shares, masterBlindingFactor: qCoeffs[0].toString(16) }
}
\end{lstlisting}

\subsubsection{同态份额刷新}

\textbf{创新核心}。传统的份额刷新需要解密主密钥 $\rightarrow$ 重新生成 Shamir 多项式 $\rightarrow$ 重新加密，整个过程依赖主密钥的可用性。而在本系统的安全模型中，主密钥在创建后即被销毁，解密重构方法不可行。

\vspace{0.5em}
\noindent \textbf{解决方案}。利用 Pedersen 承诺的同态加法性质，生成两个零和多项式 $Z_v(x)$ 和 $Z_r(x)$，满足 $Z_v(0) = Z_r(0) = 0$。由于常数项为零，主承诺不变：

\[
C'_{\text{master}} = C_{\text{master}} \quad (\text{因为 } Z_v(0) = Z_r(0) = 0)
\]

每份份额的承诺更新为：

\[
C'_i = C_i + Z_r(i) \cdot G + Z_v(i) \cdot H
\]

全程仅涉及椭圆曲线的加法和标量乘法运算，无需接触主密钥。

\begin{lstlisting}[language=TypeScript, caption=零和多项式同态刷新（shamir.ts）]
generateRefreshDeltas(n: number, t: number): ShareDelta[] {
    const zvCoeffs: bigint[] = [BigInt(0)]    // Z_v(0) = 0
    const zrCoeffs: bigint[] = [BigInt(0)]    // Z_r(0) = 0
    for (let i = 1; i < t; i++) {
        zvCoeffs.push(generateRandomCoefficient())
        zrCoeffs.push(generateRandomCoefficient())
    }
    for (let i = 1; i <= n; i++) {
        const x = BigInt(i)
        deltas.push({
            index: i,
            valueDelta: evaluatePoly(zvCoeffs, x).toString(16),
            blindingDelta: evaluatePoly(zrCoeffs, x).toString(16),
        })
    }
    return deltas
}
\end{lstlisting}

\subsubsection{Pedersen 承诺验证}

基于 secp256k1 椭圆曲线实现的 Pedersen 承诺：

\begin{lstlisting}[language=TypeScript, caption=Pedersen 承诺生成（shamir.ts）]
generatePedersenCommitment(value: string, blindingFactor?: string) {
    const H = this.getH()              // H = G * SHA256(seed)
    const G = this.ec.g                // secp256k1 生成点
    const C = G.mul(rBN).add(H.mul(vBN))  // C = r*G + v*H
    return { commitment: C.encode('hex'), blindingFactor: r }
}
\end{lstlisting}

\subsubsection{主承诺同态验证}

\begin{lstlisting}[language=TypeScript, caption=主承诺同态验证（shamir.ts）]
verifyMasterCommitment(shares, masterCommitment): boolean {
    let sum = G.mul(new BN(0))  // 从无穷远点开始
    for (let i = 0; i < t; i++) {
        // 计算 Lagrange 系数 lambda_i
        let lambda = new BN(1)
        for (let j = 0; j < t; j++) {
            if (i !== j) {
                lambda = lambda.mul(num).mul(den.invm(curveN)).mod(curveN)
            }
        }
        // sum += lambda_i * C_i （EC 点乘）
        sum = sum.add(point.mul(lambda))
    }
    return sum.encode('hex', false) === masterCommitment
}
\end{lstlisting}

\subsection{后端业务逻辑模块}

后端业务逻辑模块（legacyPlanService.ts）封装了所有核心业务流程：

\begin{itemize}
    \item \textbf{createPlan()}：接收用户输入，调用密码学引擎生成主密钥、双多项式份额和承诺，加密资产，存储数据并通过邮件分发份额，销毁主密钥。
    \item \textbf{refreshPlanShares()}：生成零和多项式增量，更新存储的承诺（EC 点运算），通过邮件发送增量给所有监护人。全程不接触主密钥。
    \item \textbf{initiateInheritance()}：验证时间锁状态，创建继承请求，通知监护人。
    \item \textbf{submitGuardianShare()}：验证份额值对 Pedersen 承诺的匹配性（支持正常和胁迫两种承诺），收集份额至门限后恢复主密钥并解密资产。
\end{itemize}

\subsection{前端功能模块}

前端提供了完整的用户交互界面\footnote{需补充功能截图}：

\begin{itemize}
    \item \textbf{Home 首页}：项目介绍和登录入口。
    \item \textbf{CreatePlan 创建计划}：表单式创建遗产计划，包含资产列表、监护人信息、门限配置和触发模式设置。
    \item \textbf{Dashboard 控制台}：展示所有计划（创建的、继承的、监护的），支持详情查看和份额刷新触发。
    \item \textbf{Inheritance 继承}：发起继承请求，输入计划 ID 和继承人信息。
    \item \textbf{Guardian 监护人}：提交份额值，查看份额状态。
    \item \textbf{AIHelper AI 助手}：基于 Deepseek 大模型 API 的自然语言对话界面，支持文件上传和上下文记忆。
\end{itemize}

\begin{figure}[H]
    \centering
    \vspace{3em}
    \framebox{\parbox{0.8\textwidth}{\centering \textbf{图 2}：系统功能界面截图 \\
    \small\textit{需补充：创建计划页面 / 控制台 / 监护份额提交页面的截图}}}
    \vspace{1em}
    \caption{系统功能界面示例}
    \label{fig:screenshots}
\end{figure}

\subsection{邮件服务模块}

邮件服务模块（emailService.ts）负责自动生成和发送各类通知邮件，包括：

\begin{itemize}
    \item \textbf{份额通知邮件}：包含份额值、盲因子、胁迫份额值及使用说明。
    \item \textbf{刷新通知邮件}：包含份额增量 $\delta_{v,i}$ 和 $\delta_{r,i}$，及新胁迫份额替换说明。
    \item \textbf{继承请求通知}：通知监护人有人发起继承请求，需确认并提交份额。
    \item \textbf{继承成功通知}：向继承人发送解密后的资产信息。
    \item \textbf{胁迫警报}：当检测到胁迫份额提交时，立即通知所有监护人。
\end{itemize}

所有邮件模板基于 HTML 构建，支持富文本显示和安全警告提示。
